\clearpage
\phantomsection% 
\addcontentsline{toc}{chapter}{RINGKASAN}
\thispagestyle{fancy}

\begin{center}
	\large \bfseries \MakeUppercase{Ringkasan}\\
	\normalsize \normalfont {\thetitle}\\
	\normalsize \normalfont {\theauthor}\\
	\bigskip
	
	\normalsize \normalfont \justifying
	\normalsize \normalfont \justifying
	% P1: Latar Belakang
	Indonesia merupakan salah satu negara dengan kekayaan sumber daya mineral yang melimpah, termasuk timah. Kasiterit (\senyawatin) merupakan mineral utama pembentuk timah yang secara teoritis mengandung 78\% timah, namun kandungan aktualnya hanya sekitar 73–75\% karena jarang ditemukan dalam bentuk murni. Proses identifikasi mineral kasiterit di Laboratorium Eksplorasi PT. Timah Tbk. saat ini masih mengandalkan metode konvensional \textit{Grain Counting Analysis} (GCA) yang dilakukan secara manual di bawah mikroskop binokuler. Metode ini memiliki sejumlah keterbatasan, di antaranya bersifat subjektif, rentan terhadap kelelahan visual (\textit{eye fatigue}), serta menghasilkan variasi hasil identifikasi antaranalis. Selain itu, keberadaan butiran mineral yang saling tumpang tindih (\textit{occluded}) menyebabkan kesulitan dalam membedakan batas antar butir kasiterit sehingga menurunkan akurasi hasil identifikasi. Oleh karena itu, untuk menangani masalah komputasional karena keterbatasan metode konvesional adalah dengan menggunakan pendekatan berbasis teknologi yang mengidentifikasi mineral pada citra mikroskop menggunakan  \textit{deep learning}. \par

	% P2: Tujuan Penelitian
	Penelitian ini bertujuan untuk: (1) menghasilkan model identifikasi berbasis \textit{Mask} R-CNN yang dapat melakukan segmentasi mineral kasiterit pada citra mikroskop serta \textit{robust} terhadap oklusi melalui teknik \textit{amodal instance segmentation}, sehingga model mampu memprediksi bentuk utuh setiap butir kasiterit meskipun sebagian tertutup; dan (2) melakukan pengujian dan analisis untuk mengetahui performa model \textit{Mask} R-CNN menggunakan metrik \textit{Average Precision} (AP) untuk menilai ketepatan klasifikasi area mineral dan \textit{Intersection over Union} (IoU) untuk mengukur akurasi tumpang tindih antara area segmentasi prediksi dengan area sebenarnya dalam mengidentifikasi mineral kasiterit dari citra mikroskop. \par

	% P3: Rumusan Masalah dan Batasan Masalah
	Berdasarkan permasalahan tersebut, penelitian ini merumuskan dua masalah utama, yaitu: (1) bagaimana merancang dan mengembangkan model identifikasi berbasis \textit{Mask} R-CNN yang dapat melakukan segmentasi mineral kasiterit pada citra mikroskop dan mengatasi oklusi antar butir melalui teknik \textit{amodal instance segmentation}, dan (2) bagaimana mengevaluasi performa model \textit{Mask} R-CNN menggunakan metrik AP dan IoU dalam mengidentifikasi mineral kasiterit dari citra mikroskop. Adapun batasan dalam penelitian ini adalah bahwa penelitian hanya berfokus pada identifikasi mineral kasiterit (\senyawatin) dan tidak mencakup identifikasi mineral ikutan lainnya. Dataset yang digunakan sebanyak 142 citra mikroskop trinokuler yang bersumber dari Laboratorium Eksplorasi PT. Timah Tbk. yang berlokasi di Kota Pangkalpinang, Bangka Belitung, dengan proses segmentasi data dilakukan menggunakan perangkat lunak LabelMe yang terintegrasi dengan \textit{Segment Anything}. \par
	
	% Pemilihan \textit{Mask} R-CNN didasarkan pada kemampuannya dalam melakukan deteksi dan \textit{instance segmentation} secara presisi pada citra berlatar belakang kompleks, sementara pendekatan \textit{amodal} memungkinkan model memprediksi bentuk utuh objek meskipun sebagian tertutup. \par

	% P4: Metodologi Penelitian (diperluas)
	Metodologi penelitian diawali dengan tahap identifikasi masalah, di mana ditemukan bahwa proses identifikasi mineral kasiterit di Laboratorium Eksplorasi PT. Timah Tbk. masih dilakukan secara manual menggunakan metode konvensional GCA. Tahap pengumpulan data menghasilkan 142 citra mikroskop mineral kasiterit yang merupakan data sekunder dari Laboratorium Eksplorasi PT. Timah Tbk. Prapemrosesan data meliputi pemilihan data berdasarkan kriteria kualitas pencahayaan dan kejelasan objek, anotasi dan segmentasi data menggunakan metode poligon yang menghasilkan 7.664 objek kasiterit, pembagian citra (\textit{image tiling}) dengan grid $2 \times 2$ yang meningkatkan jumlah objek menjadi 8.433 dengan total 537 \textit{tiles}, augmentasi data menggunakan teknik \textit{flip} horizontal dan vertikal, \textit{rotate} dengan rentang $-15$ hingga $15$ derajat, serta penambahan \textit{Gaussian noise} dengan variansi 0{,}2, dan normalisasi untuk menyeragamkan rentang nilai piksel citra. Pembagian data menggunakan metode \textit{k-fold cross validation} dengan $K=5$ dan rasio 80:20 pada setiap iterasi. Pengembangan model dilakukan dengan melatih dua model \textit{pre-trained} \textit{Mask} R-CNN berbasis \textit{backbone} ResNet50 FPN versi V1 dan V2, dengan dua jenis model \textit{standard} dan \textit{amodal} serta dua variasi \textit{optimizer} SGD dan \textit{Adam}. Evaluasi kinerja model dilakukan menggunakan metrik standar \textit{COCO} yang meliputi AP dan IoU pada beberapa ambang batas, yaitu 0{,}50, 0{,}75, dan 0{,}50:0{,}95, dengan metrik utama \APrangefiftyy dan \IoUrangefiftyy. \par

	% P5: Hasil Evaluasi dan Analisis Optimizer Adam vs SGD
	Hasil evaluasi menunjukkan bahwa model \textit{pre-trained Mask} R-CNN ResNet50 FPN V1 dengan pendekatan \textit{Amodal} dan \textit{optimizer Adam} menghasilkan performa terbaik dengan nilai rata-rata \APrangefiftyy sebesar 0{,}458 dan \IoUrangefiftyy sebesar 0{,}915. Model terbaik kedua adalah \textit{Mask} R-CNN ResNet50 FPN V2 dengan pendekatan \textit{Standard} dan \textit{optimizer} SGD yang memperoleh nilai \APrangefiftyy sebesar 0{,}449 dan \IoUrangefiftyy sebesar 0{,}915, dengan selisih \APrangefiftyy sebesar 0{,}009 terhadap model terbaik pertama. Secara keseluruhan, analisis terhadap seluruh konfigurasi pelatihan menunjukkan bahwa penggunaan \textit{optimizer Adam} secara konsisten menghasilkan performa yang lebih unggul dibandingkan \textit{optimizer} SGD. Keunggulan ini disebabkan oleh mekanisme \textit{adaptive learning rate} yang dimiliki \textit{Adam}, yang mampu mempercepat dan menstabilkan proses konvergensi terutama pada objek dengan bentuk kompleks seperti butir mineral kasiterit pada citra mikroskop. \par
	
	% Pada model V1, penggunaan \textit{Adam} menghasilkan \APrangefiftyy yang lebih tinggi baik pada pendekatan \textit{Amodal} (0{,}458 vs 0{,}446) maupun \textit{Standard} (0{,}450 vs 0{,}437), dengan selisih masing-masing sebesar 0{,}012 dan 0{,}013. Pola serupa juga terlihat pada model V2, di mana \textit{Adam} mengungguli SGD pada pendekatan \textit{Amodal} dengan selisih yang lebih signifikan sebesar 0{,}019 (0{,}420 vs 0{,}401). 
	
	% P6: Penerapan Amodal di Setiap Model Pre-trained
	Penerapan teknik \textit{Amodal Instance Segmentation} menunjukkan pengaruh yang berbeda pada masing-masing model \textit{pre-trained}. Pada model \textit{pre-trained Mask} R-CNN ResNet50 FPN V1 dengan \textit{optimizer Adam}, penambahan teknik \textit{amodal} yang meningkatkan jumlah parameter sebesar 2.950.914 (sekitar 6{,}75\%) terbukti efektif meningkatkan performa, ditunjukkan oleh kenaikan nilai \APrangefiftyy sebesar 0{,}021 (dari 0{,}437 menjadi 0{,}458) serta peningkatan \IoUrangefiftyy sebesar 0{,}005 (dari 0{,}910 menjadi 0{,}915), dengan penambahan waktu pelatihan sebesar 0{,}37 jam yang masih dalam batas komputasi yang dapat diterima. Sebaliknya, pada model \textit{pre-trained Mask} R-CNN ResNet50 FPN V2 dengan \textit{optimizer} SGD, penerapan teknik \textit{amodal} yang meningkatkan jumlah parameter sebesar 2.950.914 (sekitar 6{,}46\%) justru tidak memberikan peningkatan performa. Nilai \APrangefiftyy mengalami penurunan sebesar 0{,}029 (dari 0{,}449 menjadi 0{,}420), sementara \IoUrangefiftyy menurun sebesar 0{,}002 (dari 0{,}915 menjadi 0{,}913), meskipun waktu pelatihan hanya bertambah 0{,}45 jam. Hasil ini menunjukkan bahwa peningkatan kompleksitas model melalui pendekatan \textit{amodal} tidak selalu berbanding lurus dengan peningkatan kinerja, serta mengindikasikan bahwa efektivitas teknik \textit{Amodal Instance Segmentation} sangat bergantung pada kesesuaian arsitektur model \textit{pre-trained} dan pemilihan \textit{optimizer} yang digunakan. \par

	% P7: Hasil Pengujian dan Pengamatan Oklusi
	Pengujian kedua model terbaik dilakukan pada 6 citra uji yang berbeda dari dataset pelatihan. Hasil pengujian menunjukkan perbedaan karakteristik yang jelas antara kedua model. Model V2 \textit{Standard} dengan \textit{optimizer} SGD mampu mengidentifikasi objek kasiterit dengan akurasi deteksi yang tinggi, namun cenderung mengalami \textit{overprediction} yang mengakibatkan tingginya jumlah \textit{False Positive}. Pada data uji ke-5, model V2 \textit{Standard} bahkan menghasilkan nilai \IoUrangefiftyy berupa \textit{NaN}, yang mengindikasikan kegagalan deteksi pada ambang batas IoU tertentu. Sebaliknya, model V1 \textit{Amodal} dengan \textit{optimizer Adam} menunjukkan performa segmentasi yang lebih stabil dan konsisten, ditandai dengan nilai \IoUrangefiftyy yang lebih tinggi secara konsisten serta jumlah \textit{False Positive} yang lebih rendah, meskipun nilai \APrangefiftyy cenderung lebih rendah karena sifat prediksi yang lebih konservatif. Hasil pengamatan visual pada kondisi oklusi memperkuat temuan tersebut, di mana pada kasus kasiterit yang saling bersebelahan dan mengalami oklusi, model V1 \textit{Amodal} mampu memisahkan dan memprediksi masing-masing butir kasiterit secara individual sesuai dengan anotasi \textit{Ground Truth}, sementara model V2 \textit{Standard} cenderung memprediksi area tersebut sebagai satu objek utuh. Lebih lanjut, pada kondisi oklusi oleh elemen latar belakang berupa penanda skala (500 $\mu$m), model V1 \textit{Amodal} masih mampu mendeteksi kasiterit yang tertutup sebagian, sedangkan model V2 \textit{Standard} gagal mendeteksi dan mengabaikan objek tersebut sepenuhnya. \par

	% P8: Kesimpulan
	Berdasarkan hasil penelitian secara keseluruhan, dapat disimpulkan bahwa penelitian ini berhasil merancang dan mengembangkan model identifikasi berbasis \textit{Mask} R-CNN yang dapat melakukan segmentasi mineral kasiterit pada citra mikroskop serta mengatasi oklusi antar butir melalui teknik \textit{amodal instance segmentation}. Model \textit{pre-trained Mask} R-CNN ResNet50 FPN V1 \textit{Amodal} dengan \textit{optimizer Adam} terbukti sebagai konfigurasi terbaik dengan nilai \APrangefiftyy sebesar 0{,}458 dan \IoUrangefiftyy sebesar 0{,}915, didukung evaluasi metrik AP dan IoU melalui skema \textit{k-fold cross validation} yang memberikan gambaran komprehensif terhadap kinerja deteksi dan kualitas segmentasi. Temuan ini menegaskan bahwa pendekatan \textit{amodal instance segmentation} terbukti lebih efektif dan andal untuk identifikasi mineral kasiterit pada citra mikroskop dengan tingkat oklusi yang kompleks dibandingkan pendekatan \textit{standard}. \par

	% P9: Manfaat, Harapan, dan Saran
	Penelitian ini diharapkan dapat memberikan kontribusi sebagai fondasi bagi pengembangan model identifikasi mineral berbasis \textit{deep learning} serta menjadi referensi bagi penelitian selanjutnya yang berkaitan dengan penerapan teknologi \textit{deep learning} dalam bidang geologi dan pertambangan. Untuk pengembangan lebih lanjut, disarankan agar penelitian selanjutnya menerapkan identifikasi pada jenis mineral lain selain kasiterit, mengembangkan model segmentasi yang dapat bekerja secara \textit{real-time} langsung dari mikroskop, menerapkan strategi untuk mengurangi \textit{overprediction}, menerapkan identifikasi jenis kasiterit yang lebih spesifik berdasarkan ciri visual seperti warna, tekstur, atau bentuk kristal, serta menggunakan dua dataset terpisah yakni dataset \textit{amodal} dan \textit{modal} untuk melihat perbedaan hasil penerapan teknik \textit{Amodal Instance Segmentation}. \par


	\vfill
	
\end{center}
\clearpage